% CHU SO HUU: Ngọc Trâm | KHONG THEO PREAMBLE, KHONG BEGIN DOCUMENT TRONG THAN PHAN NAY.
\section{Ước lượng khoảng tin cậy Wilson cho tỷ lệ}
\label{sec:phuongphap-wilson}

\subsection{Mục đích và tham số}
Nhóm sử dụng khoảng tin cậy Wilson để lượng hóa độ bất định của tỷ lệ năm nhóm
Normal, DoS, Probe, R2L và U2R trong KDDTrain+. Với nhóm $g$, tham số cần ước lượng là
$p_g=P(Y=g)$; nếu có $x_g$ trong tổng số $n$ bản ghi thuộc nhóm đó thì ước lượng điểm là
$\widehat p_g=\frac{x_g}{n}$. Mỗi nhóm được xét như một biến chỉ báo nhị thức ``thuộc
nhóm $g$'' hoặc ``không thuộc nhóm $g$''.

\subsection{Công thức}
Với mức tin cậy $1-\alpha$ và $z=z_{1-\frac{\alpha}{2}}$, khoảng Wilson là
\[
\frac{
  \widehat p+\frac{z^2}{2n}
  \ \pm\ z\sqrt{\frac{\widehat p(1-\widehat p)}{n}+\frac{z^2}{4n^2}}
}{1+\frac{z^2}{n}}.
\]
Nghiên cứu dùng $\alpha=0{,}05$, do đó $z\approx1{,}96$. So với khoảng Wald
$\widehat p\pm z\sqrt{\frac{\widehat p(1-\widehat p)}{n}}$, Wilson ổn định hơn khi
tỷ lệ gần 0 hoặc 1 và không tạo cận ngoài miền $[0,1]$; đặc điểm này quan trọng đối
với nhóm U2R rất hiếm.

\subsection{Giả định và điều kiện áp dụng}
Các bản ghi được giả định độc lập, nhãn nhóm loại trừ lẫn nhau và phép ánh xạ nhãn
không sai. Khoảng Wilson không đòi hỏi xấp xỉ chuẩn của $\widehat p$ tốt như khoảng
Wald, nhưng vẫn dựa trên mô hình nhị thức và không tự xử lý sự phụ thuộc giữa các kết
nối. Năm khoảng được tính riêng, không phải một miền tin cậy đồng thời cho toàn bộ
vector $(p_1,\ldots,p_5)$; sự chồng lấp hay không chồng lấp của chúng cũng không thay
thế một kiểm định so sánh tỷ lệ.

\subsection{Cách sử dụng trong bộ dữ liệu}
Nhóm đếm năm nhãn trên toàn bộ 125.973 bản ghi KDDTrain+, không gộp KDDTest+ và không
dùng dữ liệu SMOTE. Kết quả chỉ mô tả cơ cấu của mẫu train đã tuyển chọn, không được
diễn giải thành tỷ lệ tấn công của Internet hiện tại. Ước lượng điểm, cận dưới và cận
trên được trình bày tại Mục~\ref{sec:ketqua-wilson}.

